\SourceLeaf{45}
\[ \cos a \cos b = \tfrac12\cos(a-b)+\tfrac12\cos(a+b). \]

\par\noindent\hspace*{\parindent}Man nehme nun $\cos\varphi$. Mit $\cos\varphi$ multipliziert ergibt dies $\cos\varphi . \cos\varphi$. Führt man die Multiplikation aus, so wird dieses Produkt durch $\cos\varphi^2$ dargestellt. Löst man es dagegen nach der eben angegebenen Formel in eine Summe auf, so wird es durch $\tfrac12+\tfrac12\cos 2\varphi$ dargestellt. Es ist also $\cos\varphi^2 = \tfrac12+\tfrac12\cos 2\varphi$. Geht man auf dieselbe Weise zu höheren Potenzen über, so erhält man nach und nach die folgenden Formeln

\[ \cos\varphi^3 = \tfrac34\cos\varphi + \tfrac14\cos 3\varphi ; \]
\[ \cos\varphi^4 = \tfrac38 + \tfrac48\cos 2\varphi + \tfrac18\cos 4\varphi ; \]
\[ \cos\varphi^5 = \tfrac{10}{16}\cos\varphi + \tfrac{5}{16}\cos 3\varphi + \tfrac{1}{16}\cos 5\varphi ; \]
\[ \cos\varphi^6 = \tfrac{10}{32} + \tfrac{15}{32}\cos 2\varphi + \tfrac{6}{32}\cos 4\varphi + \tfrac{1}{32}\cos 6\varphi ; \]
\[ \cos\varphi^7 = \tfrac{35}{64}\cos\varphi + \tfrac{21}{64}\cos 3\varphi + \tfrac{7}{64}\cos 5\varphi + \tfrac{1}{64}\cos 7\varphi ; \]
\[ \cos\varphi^8 = \tfrac{35}{128} + \tfrac{56}{128}\cos 2\varphi + \tfrac{28}{128}\cos 4\varphi + \tfrac{8}{128}\cos 6\varphi + \tfrac{1}{128}\cos 8\varphi . \]

\par\noindent Aus diesen Rechnungen lässt sich das allgemeine Gesetz leicht ableiten. Bei geraden Potenzen von $\cos\varphi$ treten nur Kosinus gerader Vielfacher von $\varphi$ auf, deren größtes durch den gegebenen Exponenten bestimmt ist. Bei ungeraden Potenzen treten nach derselben Regel Kosinus ungerader Vielfacher auf. Die geraden Potenzen folgen jedoch nicht ganz demselben Gesetz wie die ungeraden.

$\alpha$. Das mit $\cos 0\varphi = 1$ multiplizierte Anfangsglied bei den geraden Potenzen hat einen Koeffizienten, dessen Nenner wie bei allen übrigen Gliedern eine Potenz von 2 ist, deren Exponent um eins kleiner ist als der Exponent der vorgegebenen Potenz. Der Zähler ist im Wesentlichen ein zur gegebenen Potenz gehöriger Binomialkoeffizient, und zwar die Hälfte des mittleren, wenn man diese Koeffizienten der Reihe nach entwickelt. Die übrigen Glieder enthalten ein Vielfaches von $\varphi$, das dem Doppelten des Index des Gliedes entspricht. In ihrem Koeffizienten steht derselbe Nenner wie im Anfangsglied. Als Zähler steht darin ein Binomialkoeffizient derselben Potenz, der um so viele Stellen auf den des Anfangsgliedes folgt, wie der Index des Gliedes angibt. Dieses Gesetz lässt sich allgemein so angeben. Die Reihe für $\cos\varphi^{2m}$ hat das Anfangsglied\HistoricalNote{*)}{*)}{\textsc{Euler}s Ausdruck für dieses Gesetz, der nicht auf Binomialkoeffizienten zurückgeführt ist, stimmt der Sache nach mit dem vorgeschlagenen überein, scheint aber weniger übersichtlich.

Ich brauche kaum zu erwähnen, dass ich Binomialkoeffizienten stets nach der Schreibweise \textsc{Hindenburg}s bezeichne und, dem berühmten\ \textsc{Kaestner} folgend, bei Reihen stets das Anfangsglied von den folgenden Gliedern unterscheide.}

\[ \frac{1}{2^{2m-1}} . \frac{{}^{2m}\mathfrak{B}^{(m)}}{2}, \]
